\documentstyle[%


righttag%


,11pt%


,amssymb%


,russian%


,izvru%


]{amsart}





\newcommand{\tl}[3]{\medskip\noindent{\bf#1}\ #2 \dotfill #3 \par} 


\pagestyle{empty}


\begin{document}


%\tableofcontents


\par


\vskip 2cm


\par


\bigskip


\bigskip


\bigskip


\par


\centerline{\Large СОДЕРЖАНИЕ}


\bigskip


\bigskip


%%%%%%%%%% Use \newline %%%%%%%%%%%





\tl{Ахмедзянова Е.Г., Дубинин В.Н.}


{Радиальные преобразования множеств


и неравенства\newline для трансфинитного диаметра}{3}


\tl{Бикчантаев И.А.}


{Теорема Клиффорда для квазирациональных функций}{9}


\tl{Бобочко В.Н.}


{Краевая задача для системы сингулярно


возмущенных\newline дифференциальных уравнений


с недиагонализируемым предельным оператором}{14}


\tl{Вершинин Р.В.} 


{Об $(1+\varepsilon_n)$-ограниченных $M$-базисах}{24}


\tl{Волкодавов В.Ф., Захаров В.Н.} 


{Экстремальные свойства решений одного 


уравнения\newline гиперболического типа третьего 


порядка в трехмерном пространстве  


и их применение}{28}


\tl{Загиров Н.Ш.} 


{Полиномиальные ужи по подсистемам


алгебраических степеней}{32}


\tl{Пинус А.Г.} 


{Об элементарных теориях генерических мультиалгебр}{39}


\tl{Тептин А.Л.} 


{К вопросу об осцилляционности спектра


многоточечной краевой задачи}{44}


\tl{Тюрин С.А.} 


{Перестройки торов в алгебре Цассенхауза}{54}


\tl{Юдин М.Д.} 


{Об обобщении формулы Колмогорова на 


суммы зависимых векторов}{61}





\bigskip


{\centerline {КРАТКИЕ СООБЩЕНИЯ}}


\bigskip





\tl{Ганиев И.Г.} 


{О векторных мерах со значениями в  


пространствах Банаха-Канторовича}{65}


\tl{Здоровенко М.Ю.} 


{Гиперконечная аппроксимация преобразования Фурье}{68}


\tl{Куликова Т.Ю.}


{Приближение функций в метрике $L_2$ с весом Якоби}{73}


\tl{Лазарев Ю.В., Марченко А.В., Симоненко И.Б.} 


{О вычислении определителей\newline усеченных одномерных 


континуальных сверток с рациональными символами}{77}


\tl{Смолин Ю.Н.}


{Об одном методе получения экспоненциальной


оценки решения уравнения Вольтерра}{79}


\tl{Чупрунов A.Н.} 


{Принцип инвариантности для независимых наблюдений\newline


банаховозначного случайного процесса}{83}








\vfill


\noindent\raisebox{2pt}{\copyright} Казанский госудаpственный унивеpситет


\end{document}


